% ==============================================================================
% ================================== AUFGABE 3 ================================
% ==============================================================================

\chapter{Aufgabe 3: Folgemaßnahmen und Vorbereitung des Zweitgesprächs} 
\label{sec:aufgabe3}

\section{Unmittelbares Verhalten nach dem Gesprächsabbruch}
Der eskalierte Verlauf des Erstgesprächs und der anschließende Raumverlass durch die Kundin stellen eine akute Krisensituation in der Kundenbeziehung dar. Um das Mandat nicht dauerhaft zu gefährden, muss der Dienstleister unmittelbar nach dem Vorfall ein phasenbezogenes Deeskalations- und Krisenmanagement einleiten \cite[S.~42--44]{5}.

\subsection{Erregungskontrolle und emotionale Distanzierung}
Direkt nach dem Vorfall gilt der Grundsatz der Erregungskontrolle (Phase A des kooperativen Konfliktgespräches nach Berkel, zit. n. \cite[S.~49--50]{5}). Der Dienstleister darf keinesfalls im Affekt reagieren, etwa durch impulsive Gegenangriffe, rechtfertigende E-Mails oder scharfe Schuldzuweisungen. Eine unüberlegte Reaktion würde den Konflikt auf die Ebene des unkontrollierten Machtkampfes lenken und den drohenden Abbruch der Geschäftsbeziehung besiegeln. Es ist daher eine emotionale Abkühlphase einzuhalten, um die eigene Wahrnehmung zu versachlichen und die emotionale Stabilität wiederherzustellen.

\subsection{Reflexion mit Vorgesetztem oder Kollegen}
Da die Beschimpfungen eine erhebliche persönliche Kränkung darstellen, ist die eigene Erregungskontrolle ggf. nicht ausreichend. In diesem Fall sollte der Dienstleister die Situation mit einem Vorgesetzten oder Kollegen reflektieren (Nachbarschaftshilfe im Sinne Glasls, zit. n. \cite[S.~53--54]{5}), um seine emotionale Stabilität wiederherzustellen und Handlungsfähigkeit für das Zweitgespräch zu gewinnen.

\subsection{Innere Reorganisation im Modell des inneren Teams}
Kiefer empfiehlt zur kognitiven Verarbeitung der erlebten Kränkung die Durchführung einer inneren Teamkonferenz nach Schulz von Thun (Schulz von Thun, zit. n. \cite[S.~48--52]{3}). Der Dienstleister muss die inneren Stimmen ordnen:
\begin{itemize}
    \item Der \textit{gekränkte Experte} fordert die Verteidigung der eigenen fachlichen Ehre und die Zurückweisung der Beleidigungen.
    \item Der \textit{analytische Entwickler} verweist auf die fehlenden Dateninputs und das vertraglich vereinbarte Lastenheft.
    \item Der \textit{kundenorientierte Dienstleister} erkennt den immensen Vorstands- und Termindruck der Kundin und strebt den Erhalt des Kundenverhältnisses an \cite[S.~48--49]{3}.
\end{itemize}
Unter Führung des inneren „Oberhaupts“ wird die Entscheidung getroffen, den Fokus konsequent auf die Rettung des Projekts und die kundenorientierte Problemklärung zu richten \cite[S.~49--50]{3}.

\section{Kommunikative Erstmaßnahmen zur Kontaktaufnahme}
Nach der inneren Konsolidierung erfolgt die aktive, schriftliche Erstkontaktaufnahme. Hierbei kommt die Methode \textit{„Unerwünschtes aussprechen“} nach Glasl zur Anwendung (Glasl, zit. n. \cite[S.~48]{5}).

\subsection{Anwendung der Methode „Unerwünschtes aussprechen“}
Ziel dieser Methode ist es, bestehende Befürchtungen und Angstfantasien der Kundin (z.~B. die Sorge vor einem drohenden Rechtsstreit, weiterem Zeitverlust oder dem endgültigen Scheitern des Projekts vor dem Vorstand) entkräften zu lassen. Der Dienstleister richtet ein sachliches, wertschätzendes Anschreiben per E-Mail oder Brief an die Kundin, welches folgende Kernbotschaften enthält:
\begin{itemize}
    \item \textbf{Verzicht auf Schuldfragen:} Es werden keine gegenseitigen Vorwürfe bezüglich der vergangenen E-Mail-Korrespondenz oder der Beschimpfung formuliert.
    \item \textbf{Explizite Benennung der Grenzen:} Der Dienstleister spricht klar aus, dass ein dauerhafter Abbruch der Zusammenarbeit sowie ein Fehlschlagen der Automatisierung auf keinen Fall im Interesse des Dienstleisters liegen.
    \item \textbf{Signal der Lösungsbereitschaft:} Es wird versichert, dass der Dienstleister volle Verantwortung für die technische Anpassung übernimmt und primär an einer zügigen Fertigstellung der finanziellen Nachprüfung interessiert ist \cite[S.~47--48]{5}.
\end{itemize}

\subsection{Einladung zum zweiten Klärungsgespräch}
Gleichzeitig spricht der Dienstleister die Einladung zu einem zweiten, zeitnahen Klärungsgespräch aus. Um günstige Rahmenbedingungen zu schaffen, wird ein neutraler Ort (z.~B. ein Besprechungsraum im Büro des Dienstleisters oder ein ruhiger Konferenzraum) vorgeschlagen. Zudem wird ein zeitlicher Rahmen von mindestens zwei Stunden reserviert, um Zeitdruck und äußere Störungen (wie Telefonate oder E-Mails) konsequent auszuschließen \cite[S.~50--52]{5}.

\subsection{Reaktion auf ausbleibende Rückmeldung}
Sollte die Kundin auf die schriftliche Kontaktaufnahme nicht reagieren, sind gestufte Eskalationsmaßnahmen einzuleiten: Nach zwei bis drei Tagen erfolgt ein telefonischer Kontaktversuch über die Sekretariatsebene. Bleibt auch dieser ohne Wirkung, wird nach etwa einer Woche der Vorgesetzte der Kundin informiert, um die Dringlichkeit zu verdeutlichen und ggf. eine moderierende Rolle einzunehmen. Als letzte Stufe kann der eigene Vorgesetzte oder Projektleiter eingebunden werden, um die Kundenbeziehung auf höherer Ebene abzusichern.

\section{Systematische Vorbereitung auf das Zweitgespräch}
Eine erfolgreiche Durchführung des zweiten Treffens erfordert eine fundierte fachliche, organisatorische und gesprächsdramaturgische Vorbereitung.

\subsection{Fachliche Aufbereitung und Lösungsentwicklung}
Auf der Sachebene erarbeitet der Dienstleister im Vorfeld konkrete Handlungsoptionen, um den Übergang von einer konfrontativen Haltung zu einer konstruktiven Problemlösung zu ermöglichen \cite[S.~62]{5}.
\begin{itemize}
    \item \textbf{Lückenmatrix der Parameter:} Der Dienstleister erstellt eine übersichtliche Matrix, welche die 47 bestehenden Standardparameter den fehlenden Sonderfällen der Nachprüfung gegenüberstellt. Die Lücken werden nicht als „Versäumnis der Kundin“, sondern als „gemeinsam zu klärende Prozessanforderungen“ definiert.
    \item \textbf{Entwicklung von Ausweichszenarien (konsensbasierte Problemlösung):} Um der Kundin sofortigen Druck zu nehmen, erarbeitet der Entwickler ein Stufenmodell zur Nachbesserung. Die kritischsten Sonderfälle werden priorisiert vorgezogen, um die anstehende Abnahme vor dem Vorstand abzusichern.
    \item \textbf{Vereinfachtes Erfassungstemplate:} Für die künftige Zuarbeit wird ein vereinfachtes Eingabe-Formular vorbereitet, welches den Zeitaufwand auf Kundenseite auf ein Minimum reduziert.
\end{itemize}

\subsection{Strukturierung nach dem Modell des kooperativen Konfliktgesprächs}
Kiefer beschreibt, dass das Zweitgespräch dramaturgisch an den sechs Phasen nach Berkel ausgerichtet wird (Berkel, zit. n. \cite[S.~49--52]{5}).
\begin{enumerate}
    \item \textbf{Phase A (Erregung kontrollieren):} Der Dienstleister tritt von Beginn an ruhig, wertschätzend und kontrolliert auf.
    \item \textbf{Phase B (Vertrauen bilden):} Durch gezielte Ich-Botschaften und Selbstoffenbarung (z.~B. Einräumen eigener Kommunikationsdefizite bei den Zwischenständen) wird eine tragfähige Beziehungsebene aufgebaut.
    \item \textbf{Phase C (Offen kommunizieren):} Aktives Zuhören und Paraphrasieren der Aussagen der Kundin, um ihr zu signalisieren, dass ihre Notlage bezüglich der Vorstandsberichte verstanden wurde.
    \item \textbf{Phase D (Problem lösen):} Vorstellung des ausgearbeiteten Stufenmodells und der Lückenmatrix als Diskussionsbasis zur gemeinsamen Abstimmung.
    \item \textbf{Phase E (Vereinbarung treffen):} Verbindliche schriftliche Fixierung des neuen Zeit- und Arbeitsplans in einem gemeinsamen Ergebnisprotokoll.
    \item \textbf{Phase F (Persönlich verarbeiten):} Abschließende Reflexion des Gesprächs zur Sicherstellung einer nachhaltigen Arbeitsbeziehung.
\end{enumerate}

\subsection{Optionale Einbindung einer neutralen dritten Partei (Nachbarschaftshilfe)}
Sollte die Kundin dem Zweitgespräch skeptisch gegenüberstehen, schlägt der Dienstleister vor, eine neutrale dritte Partei einzubinden. Im Sinne der Nachbarschaftshilfe nach Glasl kann dies ein unbeteiligter Projektverantwortlicher oder Key Account Manager sein, der das Gespräch als moderierende Instanz leitet und die Einhaltung der Gesprächsregeln überwacht.

% ==============================================================================
% ================================== AUFGABE 3 ================================
% ==============================================================================